
% !TEX root = ./../main.tex

\section{Refinement Types in Software Verification}
\label{sec:state-of-the-art}

One of the key goals of Refinement Types is to do a seamless incorporation of formal methods 
into mainstream software, thanks to its automatic verification capabilities, the flexibility that it provides 
for building specifications via subtyping, and its inherent design that allows it to be integrated to 
existing languages~\cite{jhala21}. 
Here we develop the details about Refinement Types, and put forward the case for using Refinement 
types to verify \ac{COP} programs. 

%%
\subsection{The Problem of Verifying Context-Oriented Programs}
\label{sec:problem-of-verifying-cop}

\fref{sec:cop} presents the challenges of reasoning about \ac{COP} programs.
In general, the fact that the result of a partial method depends on the order of the composition and the current active layers, complicates the process of reasoning about the expected result; particularly so when dealing with the implicit scoping of the declarative activation mechanism.
Programs can get complicated very quickly when layers use partial methods from other layers, as the result of a call not only depends on the activation behavior of the partial method in question, but also on the activations of the external layers that are calling such method.
The problem is then to find a way of validating that the result behaves as expected for all possible interactions and activation configurations, but this space of possibilities is vast and may suffer from a combinatorial explosion.

There are some proposals to address this problem. A recent approach puts forward the use of Combinatorial Testing to generate a set of transitions between partial configurations, in order to detect incorrect behavior that arise from the activation of specific layers~\cite{martou24}.
A partial configuration is a set denoting the layers that could be active or inactive at a particular point in time of the execution~\cite{martou24}. Combinatorial Testing generates transitions between configurations in accordance to some heuristic.
Heuristics are used to maximize the coverage of interactions while optimizing for performance, as it is not practical to explore the complete space of all possible transitions.
The drawback of this is, however, that there may be some error cases that cannot be detected if the heuristic did not generate the particular transition that triggered the error.
Combinatorial Transition Testing~\cite{martou24} is a technique that was proposed to improve the coverage of transitions and detect more errors.

The objective of this work is aligned with these testing tools in the sense that both seek to prevent behavioral errors from happening due to inconsistent layer activations.
However the approach that is taken is substantially different, as the idea is to statically verify properties about these transitions to guarantee correct behavior without the need of exploring the space of transitions.
To achieve this goal, Refinement Types are used to define a specification of the valid transitions in the program, and enforce it in such a way that an erroneous implementation using invalid transitions is detected at compile-time.


%%
\subsection{Refinement Types}
\label{sec:refinement-types}

Type systems are often viewed as a way of constraining the set of possible values for terms, however, 
as briefly discussed in \fref{sec:types}, they are also a very powerful tool to prove and reason about 
different properties of a program. Refinement Types extend a type system with logical predicates to 
enable a richer specification of invariants for values and operations~\cite{jhala21}.
The canonical example of Refinement Types is the definition of natural numbers by refining the type of 
the integers as shown in \fref{lst:refinement-nat-example}. Here, \scode{Int} is the base type 
to be refined, and \scode{v} is a variable that represents any term adhering to the base type, which is 
then restricted by the \scode{0 <= v} predicate.

\begin{scala}[numbers=none,label=lst:refinement-nat-example, caption=Defining the Naturals with Refinement Types]
  type Nat = Int{v } 0 <= v}
\end{scala}

The types of the parameters and return values of a function can also be refined, effectively providing a 
way of specifying its preconditions and post-conditions respectively. In fact, it has been shown that 
Refinement Types have a deep connection with Hoare Logic~\cite{chen22}. Refinement Types also 
allow the developer to declare Dependent Function Types which are, in practical terms, functions 
whose return type may depend on the input~\cite{rijke22}. Using these features, one may describe very 
precise contracts that get checked at compile-time. For example, it is possible to declare a function for 
indexing an array while guaranteeing that no out of bounds error will occur. \fref{lst:safe-array-access} 
shows an example of this; note how the type of the \scode{index} parameter depends on the previous 
\scode{arr} parameter.

\begin{scala}[numbers=none,label=lst:safe-array-access, caption=Safe Array Access]
  def safeGet[A](arr: Array[A])(index: Nat{v} v < arr.size}): A
\end{scala}

To achieve automatic verification, Refinement Types use an \ac{SMT} solver to determine whether the 
logical constraint reflected by a type is true or false. An \ac{SMT} solver is a collection of algorithms 
that solve the decision problem for logical formulas drawn from combinations of certain theories such 
as arithmetic, arrays, or uninterpreted functions~\cite{demoura08}.
The types get translated into \ac{VC} constraints that the solver can understand to produce a result.
A \ac{VC} constraint is either a quantifier-free predicate, a conjunction of constraints, or an implication 
from a predicate to a constraint~\cite{jhala21}. Equations \ref{eq:var-eq} through \ref{eq:arith-eq} 
illustrate how the predicates and \ac{VC} constraints can be constructed. It is important that the 
predicates are decidable logic formulas so that type checking stays deterministic~\cite{jhala21}.
Even though solvers like Z3 have heuristics to work with logic theories for which there is no decision 
procedure (for instance, first-order logic)~\cite{demoura08}, troubleshooting a type checking failure 
may be harder if the types rely on such heuristics, whose implementation may vary from solver to 
solver~\cite{jhala21}.

\begin{equation}
\label{eq:var-eq}
\text{Variables } v ::= x,y,z,\dots
\end{equation}

\begin{minipage}[t]{0.5\linewidth}
\begin{equation}
\begin{split}
\text{\ac{VC} constraint } c &::= p \\
  &| \quad c \wedge c \\
  &| \quad \forall x : b \: . \: p \Rightarrow c
\end{split}
\end{equation}
\end{minipage}
\hfill
\begin{minipage}[t]{0.5\linewidth}
\begin{equation}
\label{eq:pred-eq}
\begin{split}
\text{Predicate } p &::= true \; | \; false \\
  &| \quad v \\
  &| \quad p \wedge p \\
  &| \quad p \vee p \\
  &| \quad \neg \; p \\
  &| \quad p \Rightarrow p \\
  &| \quad p \Leftrightarrow p \\
  &| \quad v(p_1, \; p_2, \; \dots, \; p_n) \\
  &| \quad a \; r \; a \\
\end{split}
\end{equation}
\end{minipage}

\begin{minipage}[t]{0.5\linewidth}
  \begin{equation}
  \begin{split}
  \text{Arithmetic Relation } r &::= \quad = \\
    &| \quad \neq \\
    &| \quad < \\
    &| \quad \le  \\
    &| \quad > \\
    &| \quad \ge \\
  \end{split}
  \end{equation}
\end{minipage}
\hfill
\begin{minipage}[t]{0.5\linewidth}
  \begin{equation}
  \label{eq:arith-eq}
  \begin{split}
  \text{Arithmetic Expression } a &::= -1,0,1,2,\dots \\
    &| \quad v \\
    &| \quad a + a \\
    &| \quad a - a  \\
    &| \quad a * a \\
    &| \quad \text{if } p \text{ then } a \text{ else } a \\
    &| \quad v(a_1, \; a_2, \; \dots, \; a_n) \\
  \end{split}
  \end{equation}
\end{minipage}
\vspace{1em}

Since having a decidable type checking is desirable, the language of the predicates for refining types 
is significantly restricted (\fref{eq:pred-eq}). Despite this, Refinement Types are capable of expressing 
complex specifications beyond basic arithmetic, due to the usage of \emph{uninterpreted functions} as 
shown in Equations \ref{eq:pred-eq} and \ref{eq:arith-eq}. Uninterpreted functions are functions for 
which the solver does not have any information other than the fact that two function applications are 
the same if the arguments are equal. This is encoded in the congruence axiom of the solver:

\begin{equation}
\forall x,y \; . \; x = y \Rightarrow f(x) = f(y)
\end{equation}

Uninterpreted functions also enable a mechanism for expressing arbitrary user-defined functions and 
polymorphic variables~\cite{jhala21}.

To illustrate how Type Checking works with Refinement Types, we present an example from 
the~\citet{jhala21} tutorial. What follows is a definition of the simply typed lambda calculus with 
Refinement Types:

\begin{equation}
\begin{split}
\text{Basic Types } b &::= int \; | \; unit \\
\text{Type Refinements } r &::= \{v \; | \; p \} \\
\text{Types } t &::= b\{r\} \; | \; x:t \to t \\
\text{Kinds } k &::= B \; | \; \star \\
\text{Typing Context } \Gamma &::= \varnothing \; | \; \Gamma;x:t  \\
\text{Terms } e &::= a \\
  &| \quad \text{let } x = e \text{ in } e \\
  &| \quad \lambda x . \; e \\
  &| \quad e(e) \\
  &| \quad e:t \\
\end{split}
\end{equation}

Note that the predicates $p$ that refine the types are drawn from the decidable formulas defined in 
\fref{eq:pred-eq}, and that there is a notion of $Kind$ within the language. A $Kind$ is a special 
\emph{sort} that we use to organize all the types in the language. To avoid the logic paradox that 
would arise from asserting that $Type : Type$, we can instead state that $Type : Kind$~\cite{ragde20}.
Particularly, we have two kind levels in this language:
\begin{enumerate*}[label=(\arabic*)]
  \item Kind $B$ is inhabited by all basic types that can be refined (\eg $int : B$) and,
  \item Kind $\star$ is inhabited by all types in the language, including those that inhabit kind $B$ (\eg $(x:s \rightarrow t) : \star$).
\end{enumerate*}

To construct the terms of this language, we must follow the inference rules defined for the type system.
As a starting point, we have the common rules from Dependent Type Theory to create functions as 
lambda terms ($\lambda$), apply functions to a given input ($\lambda$-app), and to evaluate and 
unfold global definitions ($\delta$).

\begin{multicols}{2}
\centering
  \begin{prooftree}
    \hypo{ \Gamma; x : s \vdash y : t(x)  }
    \infer1[$\lambda$]{ \Gamma \vdash \lambda x . y : (x : s) \rightarrow t(x)  }
  \end{prooftree}
  
  \begin{prooftree}
    \hypo{ \Gamma \vdash f : (x : s) \rightarrow t(x)  } \hypo{ \Gamma \vdash x : s }
    \infer2[$\lambda$-app]{ \Gamma \vdash f(x) : t(x) }
  \end{prooftree}
  
  \begin{prooftree}
    \hypo{ \Gamma \vdash t : \star }
    \infer1[$\delta$]{ \Gamma; x : t \vdash x : t }
  \end{prooftree}
\end{multicols}

Refinement Types extend this set of rules to introduce predicate refinements for the basic types.
All types and their refinements must be first organized in terms of the Kind system that presented 
previously. To achieve this, the type system has well-formedness rules that state how Kinds are 
inhabited~\cite{jhala21}

\begin{multicols}{2}
\centering
  \begin{prooftree}
    \hypo{ \Gamma; x : b \vdash p }
    \infer1[Wf-base]{ \Gamma \vdash b\{x \; | \; p\} : B }
  \end{prooftree}
  
  \begin{prooftree}
    \hypo{ \Gamma \vdash s : k_s  } \hypo{ \Gamma; x : s \vdash t : k_t }
    \infer2[Wf-fun]{ \Gamma \vdash x : s \to t : \star }
  \end{prooftree}
  
  \begin{prooftree}
    \hypo{ \Gamma \vdash t : B  }
    \infer1[Wf-kind]{ \Gamma \vdash t : \star }
  \end{prooftree}
\end{multicols}

To check that a type refinement is valid, we must build \acp{VC} from the predicates and consult 
with the \ac{SMT} solver if they are satisfiable. In order to do so, we introduce the ntailment 
rules~\cite{jhala21}

\begin{multicols}{2}
\centering
  \begin{prooftree}
    \hypo{ \text{Smt}(c) }
    \infer1[Ent-emp]{ \vdash c }
  \end{prooftree}

  \begin{prooftree}
    \hypo{ \Gamma \vdash \forall x : b \: . \: p \Rightarrow c }
    \infer1[Ent-ext]{ \Gamma; x : b \{x \; | \; p \} \vdash c }
  \end{prooftree}
\end{multicols}

The Ent-emp rule states that any \ac{VC} constraint $c$ is valid if the \ac{SMT} solver determined 
that $c$ holds. Ent-ext specifies that constraint $c$ can be derived from some context $\Gamma$ 
extended with the refinement of type $b$ for term $x$, $ x : b \{x \; | \; p \}$, if the constraint 
$\forall x : b \: . \: p \Rightarrow c$ can be derived from $\Gamma$. The Type Checker will use Ent-ext 
to generate the \ac{VC} constraints from the Refined Type, and then it will use the Ent-emp rule to 
decide if the Refinement Type is valid by consulting the \ac{SMT} solver.

Additionally, there is a notion of Subtyping that establishes relationships between the predicates of 
two Refined Types:

\begin{multicols}{2}
\centering
  \begin{prooftree}
    \hypo{ \Gamma \vdash \forall v1 : b \: . \: p1 \Rightarrow p2[v2 := v1] }
    \infer1[Sub-base]{ \Gamma \vdash b\{ v1 \: | \: p1 \} \prec : b\{ v2 \: | \: p2 \} }
  \end{prooftree}
  
  \begin{prooftree}
    \hypo{ \Gamma \vdash x : s } \hypo{ \Gamma \vdash s \prec : t }
    \infer2[Sub-subst]{ \Gamma \vdash x : t }
  \end{prooftree}
\end{multicols}

Sub-base asserts that the type refinement $b\{ v1 \: | \: p1 \}$ is a subtype of $b\{ v2 \: | \: p2 \}$, if 
$p1 \Rightarrow p2$ for all values of type $b$, while Sub-subst allows us to treat $x : s$ as a term of 
type $t$, if we can prove that $s$ is a subtype of $t$. These two rules allow us to, for instance, use a 
term of type $int\{v \: | \: 10 \leq v\}$ whenever we require a term of type $int\{v \: | \: 0 \leq v\}$.

Finally there are some primitive values and functions that we can use in our language.
These primitives have the most precise semantics encoded in their type~\cite{jhala21}.
In Equation \ref{eq:prim}, $prim$ is an assignment from a primitive value to its type.

\begin{equation}
\label{eq:prim}
\begin{split}
  prim(0) &\doteq int\{ v \: | \: v = 0 \} \\
  prim(1) &\doteq int\{ v \: | \: v = 1 \} \\
  prim(add) &\doteq (x:Int) \rightarrow (y:int) \rightarrow int\{ v \: | \: v = x + y \} \\
  prim(sub) &\doteq (x:Int) \rightarrow (y:int) \rightarrow int\{ v \: | \: v = x - y \}
\end{split}
\end{equation}

We can use this assignment as a rule to check the types of primitive values:

\begin{center}
\begin{prooftree}
    \hypo{ prim(const) \doteq A }
    \infer1[Prim]{ \vdash const : A }
\end{prooftree}
\end{center}

There are now enough rules to verify simple programs with the Refinement Type system. To show 
these rules in action, we can type check the following program (\ie build the proof tree for the given 
terms).

\begin{equation*}
  \begin{split}
    inc &: (x: int) \rightarrow int\{v = x + 1\} \\
    inc &= \lambda x . add(x)(1) \\
    fun &: int\{0 \leq v\} \rightarrow int\{0 \leq v\} \\
    fun &= \lambda x . inc(x)
  \end{split}
\end{equation*}

The $inc$ term is a function that adds 1 to its input, and the $fun$ function states that, for all integers 
greater or equal than 0, we are guaranteed to return an integer greater or equal than 0 by applying the 
$inc$ function to its input. The proof tree for the $inc$ function is the following:

\begin{center}
  \begin{prooftree}
    \hypo{ prim(add) \doteq (x:int) \rightarrow (y:int) \rightarrow int\{v = x + y\} }
    \infer1[Prim]{ x : int \vdash add : (x:int) \rightarrow (y:int) \rightarrow int\{v = x + y\} }
    \hypo{\vdash int : \star}
    \infer1[$\delta$]{x : int \vdash x : int}
    \infer2[$\lambda$-app]{ x : int \vdash add(x) : (y:int) \rightarrow int\{v = x + 1\} }
    \hypo{ prim(1) \doteq int\{v = 1\} }
    \infer1{ x : int \vdash 1 : int }
    \infer2[$\lambda$-app]{ x : int \vdash add(x)(1) : int\{v = x + 1\} }
    \infer1[$\lambda$]{ \vdash \lambda x . add(x)(1) : (x:int) \rightarrow int\{v = x + 1\} }
  \end{prooftree}
\end{center}

The proof tree simply applies the function introduction ($\lambda$) and function application 
($\lambda$-app) rules successively until reaching the primitive rules for the constant terms (1 and 
$add$), concluding that $y$ can be replaced by 1 to form the type $(x:int) \rightarrow int\{v = x + 1\}$.
This derivation can now be used to type check the $fun$ term, as follows.

\begin{center}
  \begin{prooftree}
    \hypo{ \vdots }
    \infer1{ x:int\{0 \leq v\} \vdash inc(x) : int\{v = x + 1\} }
    \hypo{ Smt(\forall x : int \: . \: 0 \leq x \Rightarrow \forall v : int \: . \: v = x + 1 \Rightarrow 0 \leq v) }
    \infer1[Ent-emp]{ \vdash \forall x : int \: . \: 0 \leq x \Rightarrow \forall v : int \: . \: v = x + 1 \Rightarrow 0 \leq v }
    \infer1[Ent-ext]{ x:int\{0 \leq v\} \vdash \forall v : int \: . \: v = x + 1 \Rightarrow 0 \leq v }
    \infer1[Sub-base]{ x:int\{0 \leq v\} \vdash int\{v = x + 1\} \prec : int\{0 \leq v\} }
    \infer2[Sub-subst]{ x:int\{0 \leq v\} \vdash inc(x) : int\{0 \leq v\} }
    \infer1[$\lambda$]{ \vdash \lambda x . inc(x) : (x:int\{0 \leq v\}) \rightarrow int\{0 \leq v\} }
  \end{prooftree}
\end{center}

Knowing that $inc(x)$ has the type $int\{v = x + 1\}$, we need to conclude that it also has the type 
$int\{0 \leq v\}$. We use the Subtyping rules to prove that $inc(x)$ has the type $int\{0 \leq v\}$ by 
generating the \ac{VC} constraint: 
$\forall x : int \: . \: 0 \leq x \Rightarrow \forall v : int \: . \: v = x + 1 \Rightarrow 0 \leq v$; which is 
deemed true by the \ac{SMT} solver. Note how the Ent-ext rule was used to construct the final \ac{VC} 
constraint, and how Ent-emp was used to call the \ac{SMT} solver.

More complicated Refinement Type languages extend this rule set to enable more powerful capabilities 
like path-sensitive reasoning for branching expressions, polymorphic types, termination checking and, 
proving theorems through the Curry-Howard interpretation of types~\cite{jhala21}.
Refinement Type based verifiers like LiquidHaskell~\cite{chen22} or Stainless~\cite{hamza19} have all 
of these features, and can be used as automated theorem provers. One can describe the properties of 
user-defined functions following the Propositions-as-Types paradigm, by stating the property as a 
proposition in the Refinement Type, and constructing a term that inhabits such type as the proof of such 
proposition. For instance, one could prove that the concatenation operation on lists ($+$) is associative:

\begin{equation*}
  listAssoc : (l1 \; l2 \; l3 : list(a)) \rightarrow unit\{ \; l1 + (l2 + l3) = (l1 + l2) + l3 \; \}
\end{equation*}

Sometimes, the SMT solver will be able to automatically prove the given proposition, but as the 
properties get more complex, the programer can also explicitly provide the proof by building a term of 
the specified type, often using a DSL defined by the same Refinement Type language.

There have been several usages of Refinement Types for developing reliable systems. The work done 
by~\citet{kuper22} on verifying distributed systems with LiquidHaskell is a good example of the 
applicability of Refinement Types in software development. In that work, Refinement Types are 
used to formally verify that CRDT implementations are correct, by proving that the application of 
operations to update the CRDT state is commutative. A CRDT was defined as a Typeclass containing 
the type definitions of the operation and the proof obligation of it being commutative~\cite{kuper22}.
Implementors of a CRDT instance are compelled to write the commutativity proof, otherwise the 
program will not compile. Having the verification system embedded within the same language that the 
programs are written on, is one of the main advantages highlighted by the proponents of Refinement 
Type Systems~\cite{jhala21}. Unlike many software verification tools that formally reason about a 
model that is separate from executable code, Refinement Types allow the developer to verify the 
actual code that gets executed directly.


%%
\subsection{The Stainless Verifier}
\label{sec:stainless}

Stainless~\cite{hamza19} is refinement type-based verifier written in Scala used to reason about Scala 
programs. We use Stainless in this work for the verification of COP programs. In particular, Stainless is 
based on System FR, which is a calculus that supports System F polymorphism, Refinement Types, 
and recursive types~\cite{hamza19}. Stainless is able to verify a purely functional subset of Scala that 
is SMT decidable. Moreover, Stainless is also capable of verifying an imperative fragment of the 
language which uses mutable objects and loops~\cite{hamza19}. The following example is a verified 
specification of the factorial function with Stainless, taken from the framework's documentation.

\begin{scala}
  def fact(n: BigInt): BigInt = {
    require(n >= 0)
    if n == 0 then BigInt(1) else n * fact(n -1)
  }.ensuring(res => res >= 0)
\end{scala}

The \scode{require} and \scode{ensuring} functions are used to describe the pre-conditions and 
post-conditions of the \scode{fact} function respectively. Stainless translates these calls to type 
refinements similar to \scode{BigInt{v : v >= 0}} in both cases. These refinements are then type 
checked as explained in \fref{sec:refinement-types}, by the generation of \ac{VC} constraints that are 
fed to an \ac{SMT} solver. Stainless also supports a termination checker based on sized 
types~\cite{hamza19}, and it is also able to prove theorems like the right-identity law of the Monoid 
formed by the concatenation (\scode{++}) operation on Lists.

\begin{scala}
  def rightUnit[A](xs: List[A]): Boolean = {
    xs ++ Nil == xs
  }.holds
\end{scala}

Though propositions in Stainless are expressed as boolean expressions and not as types, it is, in 
principle, the same mechanism of proving propositions through the Curry-Howard interpretation   
discussed in \fref{sec:refinement-types}. Additionally, Stainless supports type refinements for 
user-defined \acp{ADT}, and the definition of Typeclasses that require a proof obligation of its 
properties. The following code shows an example \scode{Semigroup} Typeclass defined in Stainless.

\begin{scala}
trait Semigroup[A]:
  def plus(x: A, y: A): A

  @law
  def assocLaw(x: A, y: A, z: A) =
    plus(x, plus(y, z)) == plus(plus(x, y), z)
\end{scala}

Any instance of \scode{Semigroup} should implement the \scode{plus} abstract member and provide 
the proof for the \scode{assocLaw} member.


\endinput

